\paragraph{与传播率控制模型的区别。}
为便于比较，将 \citet[式(1)]{MicloSpiroWeibull2022}
中的状态和恢复率分别改记为 $s,i,\gamma$，传播控制仍记为 $b$。
该系统为 $\dot s=-bsi$、$\dot i=(bs-\gamma)i$，
故完全抑制 $b=0$ 时，易感比例保持不变而感染比例按恢复率下降。
相比之下，在本文常接触率情形，$q=1$ 时
\[
\dot s=-csi<0,\qquad \dot i=-\gamma i<0.
\]
因此，完全跟踪不仅使感染生成项为零，还使未隔离易感者继续移出。
相应轨道具有式\eqref{eq:Psi}的不变量，
其最优切换位置由式\eqref{eq:switch-condition}的特征线匹配确定。
本文证明的非平凡切换曲线满足式\eqref{eq:Gamma-slope}，
并在出现正长度容量阶段时于 $s_B>h$ 离开容量线、进入完全跟踪阶段；
最优结构及其初值判据集中见推论\ref{cor:optimal-synthesis}。
这些是本模型相对于上述标准传播率控制模型的结构差异。

两问题的控制集和成本亦不同。
\citet[式(2)]{MicloSpiroWeibull2022}允许传播率超过其自然水平 $\beta$，
且运行成本 $[\beta-b]_+$ 在 $b\ge\beta$ 时为零；
本文则在 $q\in[0,1]$ 上使用 $pcq$，其唯一零点为 $q=0$。
因此，不应直接移用文献中关于最优控制集合或后期唯一性的结论。
本文分别验证全局最优性，并分析验证等号条件下的可行轨道唯一性。
